\chapter{Мішок із деталями}
% full version: chapters/01b_neurontypes.tex

\pencilfig{1}{\pencilart{1}}{Мішок із деталями: майже всі сині, трохи червоних і жменька решти.}

Уявіть, що вам дали мішок, у якому лежить вісімдесят шість мільярдів однакових на вигляд деталей, і попросили з'ясувати, як із них зібрано комп'ютер. З чого б ви почали? Інженер не став би роздивлятися форму деталей. Він запитав би: скільки в деталі входів, скільки виходів і що вона видає на вихід.

Нейрон --- так називається така деталь --- влаштований напрочуд просто. Входів у нього тисячі: до нього підходять проводи від інших нейронів. Вихід один. Але цей вихід галузиться, як дерево, і та сама копія сигналу йде до тисяч отримувачів. Уявіть людину, яка слухає тисячу співрозмовників одразу, а відповідає однією фразою, але так, що її чують ще тисяча людей.

Що це за сигнал? Коротке клацання. Нейрон час від часу видає імпульс --- сплеск напруги завдовжки близько тисячної частки секунди. Усі клацання одного нейрона однакові, як удари одного барабана. Тому все, що нейрон може сказати, записано не в гучності клацання, а в тому, \emph{коли} він клацає: частіше, рідше, одразу чи із затримкою.

Усередині нейрон робить найпростіше обчислення, яке тільки можна вигадати: додає вхідні сигнали і, якщо сума перевалила за поріг, клацає сам. Одні входи штовхають суму вгору, інші --- вниз.

\section*{Як розсортувати мішок}

Кінчик вихідного проводу, дотягнувшись до сусіда, не торкається його. Між ними лишається вузька щілина, і сигнал через неї передає хімія: провід викидає дрібку особливої речовини --- медіатора, --- а сусід її ловить. І ось удача: майже кожен нейрон на всіх кінчиках свого проводу викидає одну й ту саму речовину. Отже, нейрони можна сортувати за тим, \emph{що} вони викидають.

Тип нейрона позначатимемо чотирма латинськими літерами --- початком міжнародної назви його речовини. Нейрон, що викидає глутамат (glutamate), --- \nt{GLUT}. Той, що викидає гамма-аміномасляну кислоту (GABA), --- \nt{GABA}. Далі будуть \nt{DOPA} (дофамін), \nt{SERO} (серотонін), \nt{NORA} (норадреналін), \nt{ACET} (ацетилхолін) і ще кілька.

Коли мішок розсортовано, картина виходить дуже нерівна. З кожної тисячі нейронів близько дев'ятисот шістдесяти --- це \nt{GLUT}. Їхній сигнал підштовхує отримувача до клацання: це «плюс». Ще близько тридцяти шести --- \nt{GABA}, їхній сигнал відштовхує від порогу: це «мінус». У корі їх більше, від п'ятої частини до чверті. А всі інші типи разом --- кілька нейронів на сто тисяч. Крапля в морі.

\section*{Телефон і радіо}

Але ця крапля влаштована зовсім інакше. Провід \nt{GLUT}- чи \nt{GABA}-нейрона веде до певних сусідів, як телефонна лінія: сигнал отримують лише ті, кому дзвонять. А крихітні групи \nt{DOPA}, \nt{SERO}, \nt{NORA} і \nt{ACET} працюють як радіостанції. Їхні проводи розходяться величезними ділянками мозку, речовина часто виливається не в щілину, а в простір навколо, і те саме повідомлення чують одразу мільйони клітин.

Телефонними лініями йдуть дані: ось лінія на картинці, ось звук певної висоти. Радіо передає загальні налаштування: наскільки все гучно, чи варто вчитися, чи час спати. У будь-якому комп'ютері так само: комірок пам'яті мільярди, а регістрів керування --- жменька. Без них машина однаково не працює.

\section*{Зміст обирає отримувач}

Є одна тонкість. Сама речовина не знає, «плюс» вона чи «мінус». Молекула --- це ключ. Що станеться, коли ключ вставлять, вирішує замок --- особливий білок на боці отримувача. Той самий дофамін одних отримувачів підбадьорює, а інших приглушує, залежно від того, який замок у них стоїть. Це як та сама радіопередача, яку різні слухачі розуміють по-різному.

\section*{Сигнал «краще, ніж очікувалося»}

Найвідоміший радіоканал --- дофамін. Що він передає? Ось дослід. Тварині несподівано дають краплю соку, і дофамінові нейрони відповідають спалахом клацань. Потім перед соком щоразу вмикають лампочку. Тварина засвоює, що лампочка обіцяє сік, і тепер спалах буває на лампочку, а на сам сік --- ні. А якщо після лампочки соку не дають, дофамінові нейрони в момент очікування \emph{замовкають}.

\pencilfig{101}{%
% three rows: a spike raster of a DOPA neuron; lamp (amber) and juice (blue drop) above the axis
\foreach \row/\y in {1/0,2/-1.05,3/-2.1}{
  \draw[pen] (0,\y) -- (6.2,\y);
  \foreach \s in {0.3,0.75,1.1,2.15,2.6,3.05,3.45,3.9,4.45,4.95,5.4,5.85}{
    \pgfmathtruncatemacro{\gap}{(\row==3)&&(\s>3.6)&&(\s<4.7)}
    \ifnum\gap=0 \draw[pcGray, line width=0.7pt] (\s,\y) -- (\s,\y+0.3);\fi}}
% row 1: juice without a lamp -> burst at the juice
\foreach \s in {4.0,4.08,4.16,4.24,4.33}{\draw[pcAmber, line width=0.8pt] (\s,0) -- (\s,0.45);}
\draw[dotpen=pcBlue, wash=pcBlue] (4.15,0.72) circle (0.09); \draw[pcBlue, line width=0.6pt] (4.07,0.76) -- (4.15,0.92) -- (4.23,0.76);
% row 2: lamp -> burst at the lamp, nothing at the promised juice
\foreach \y in {-1.05,-2.1}{
  \foreach \s in {1.5,1.58,1.66,1.74,1.83}{\draw[pcAmber, line width=0.8pt] (\s,\y) -- (\s,\y+0.45);}
  \draw[dotpen=pcAmber, wash=pcAmber] (1.65,\y+0.72) circle (0.1);
  \foreach \a in {30,90,150}{\draw[pcAmber, line width=0.5pt] ($(1.65,\y+0.72)+(\a:0.14)$) -- ($(1.65,\y+0.72)+(\a:0.24)$);}}
\draw[dotpen=pcBlue, wash=pcBlue] (4.15,-0.33) circle (0.09); \draw[pcBlue, line width=0.6pt] (4.07,-0.29) -- (4.15,-0.13) -- (4.23,-0.29);
% row 3: lamp, no juice -> a gap where the juice was expected
\draw[pcBlue, line width=0.6pt, densely dotted] (4.15,-1.38) circle (0.09); \draw[pcBlue, line width=0.6pt, densely dotted] (4.07,-1.34) -- (4.15,-1.18) -- (4.23,-1.34);
\draw[penlite=pcRed] (3.65,-2.22) -- (3.65,-2.28) -- (4.65,-2.28) -- (4.65,-2.22);
\node[lbl, text=pcRed, anchor=north] at (4.15,-2.3) {провал};
\node[lbl, anchor=east, align=right] at (-0.1,0.15) {сік\\без лампочки};
\node[lbl, anchor=east, align=right] at (-0.1,-0.9) {лампочка,\\потім сік};
\node[lbl, anchor=east, align=right] at (-0.1,-1.95) {лампочка,\\соку немає};
}{Спалах на несподіваний сік; на лампочку, яка його обіцяє; і провал тієї миті, коли обіцяний сік не надійшов.}

Правило, що пояснює всі три випадки: дофамін повідомляє не «добре», а «краще, ніж я очікував». Несподіваний сік --- приємний сюрприз. Обіцяний --- жодного сюрпризу. Обіцяний, але не даний --- неприємний. Саме таку величину використовують програми, що вчаться на своїх помилках. До неї ми ще повернемося.

\section*{Що взяти з собою}

Мозок --- це величезна мережа однакових на вигляд деталей. Майже всі вони --- «плюси» й «мінуси» телефонної мережі даних. І зовсім небагато --- радіостанції, які налаштовують усю мережу одразу. У наступному розділі ми перевіримо, що можна побудувати з самих лише «плюсів».

\fullversion{1}
